\begin{itemize}
\item Walk-forward evaluation is materially harder than a single split because it exposes the model to regime changes and repeated re-training, which typically reduces stability.
\item Comparing Step 2 (leaky scaling) to Step 3 (fold-specific scaling + purging + embargo) shows how leakage can inflate backtest performance; the Step 3 results are the preferred estimate of deployable performance.
\item Under leakage-reduced testing, the strongest configuration in this run is \textbf{lstm} in \textbf{Step 3b} with Sharpe=1.253, TotalRet=0.996, MaxDD=-0.192.
\item Practical deployment should add position sizing, transaction-cost stress tests, and a strict research protocol (no tuning on test folds) to reduce the probability of backtest overfitting (L\'opez de Prado; Bailey et al.).
\item In broader asset-pricing contexts, neural networks and other ML methods can add economic value, but only under careful out-of-sample evaluation (Gu, Kelly, and Xiu).
\item If performance is only marginal after leakage reduction, prioritize robustness (simpler models, fewer degrees of freedom, and longer test windows) over headline returns.
\end{itemize}
